\documentclass[12pt,a4paper]{article}
\usepackage[utf8]{vietnam}
\usepackage{amsmath,amssymb,tikz,tkz-tab,multirow,booktabs,xcolor,enumitem}
\usepackage[top=1.5cm,bottom=1.5cm,left=2cm,right=1.5cm]{geometry}
\newcommand{\diem}[1]{\textcolor{blue}{(#1 điểm)}}
\newcommand{\vecto}[1]{\overrightarrow{#1}}
\def\tapR{\mathbb{R}}
\DeclareMathOperator{\Ln}{Ln}
\newtheorem{dl}{Định lí}
\newenvironment{khunglg}{\par\noindent\textbf{Hướng dẫn:}}{\par}
\newif\ifdapan \dapantrue
\title{ĐỀ KIỂM TRA GIỮA KÌ I}
\author{Tổ Toán -- Trường THPT ABC}
\date{}
\begin{document}
\maketitle
\section*{PHẦN I. Trắc nghiệm}


\textbf{Câu 1.}  Cho tập hợp $A=\{x\in\mathbb{R} \mid x^2-3x+2=0\}$. Số phần tử của $A$ là (0,25 điểm)
\boncot{$2$.}{$3$.}{$5$.}{$4$.}



\textbf{Câu 2.}  Giá trị của $\lim\limits_{x\to 0}\dfrac{\sin 3x}{x}$ bằng
\boncot{$9$.}{$5$.}{$7$.}{$11$.}



\textbf{Câu 1.} %[0D1Y1-1]
Trong mặt phẳng $Oxy$, cho $\overrightarrow{AB}=(1;2)$ và $\vv{AC}=(3;-1)$. Tính $\left\lvert\overrightarrow{AB}\cdot\vv{AC}\right\rvert$.
\choice{$1$}{\True $2$}{$3$}{$4$}


\textbf{Lời giải.} Ta có $\overrightarrow{AB}\cdot\vv{AC}=1\cdot 3+2\cdot(-1)=1$, suy ra \textbf{đáp án} là $1$.








\textbf{Câu 2.} 
Mệnh đề nào sau đây \textbf{đúng}?
\choiceTF{Hàm số $y=x^3$ đồng biến trên $\mathbb{R}$}{Đ}{Phương trình $x^2+1=0$ có nghiệm}{S}{$\sqrt{4}=2$}{Đ}{$\pi$ là số hữu tỉ}{S}




\section{Tự luận}


\textbf{Định lí 1.} (Pytago) 
Tam giác $ABC$ vuông tại $A$ khi và chỉ khi $BC^2=AB^2+AC^2$.






a) Giải phương trình $\heva{&x+y=3\\&x-y=1}$.

b) Tính tích phân
\begin{align}
I &= \int_0^1 x e^{x}\,\mathrm{d}x \label{eq:I}\\
  &= \Big[x e^x\Big]_0^1 - \int_0^1 e^x\,dx = 1. \nonumber
\end{align}
Theo ([eq:I]) ta có kết quả.

c) Chứng minh $\operatorname{Ln}(xy)=\operatorname{Ln} x+\operatorname{Ln} y$ với $x,y\in\mathbb{R}^{+}$ và $a\not\in B$, $\overparen{AB}$.





• Văn bản có chú thích (Đây là chú thích chân trang với $x^2$.) và \emph{nhấn mạnh}.

• Màu chữ đỏ và nền vàng, \underline{gạch chân}, gạch ngang.



\begin{center}
\begin{tabular}{|c|c|c|}
\hline
\multirow{2}{*}{Nhóm} & \multicolumn{2}{c|}{Kết quả} \\
\cline{2-3}
 & Đạt & Không đạt \\
\hline
A & $12$ & $3$ \\
\hline
\end{tabular}
\end{center}

\begin{table}[h]
\centering


\begin{center}\textbf{Bảng 1.} Bảng số liệu\end{center}


\begin{tabular}{lcr}
\toprule
\textbf{Tên} & Tuổi & Điểm \\
\midrule
An & 16 & 8,5 \\
Bình & 17 & 9,0 \\
\bottomrule
\end{tabular}
\end{table}
Xem Bảng~[tab:1].


 \par\noindent\textbf{Hướng dẫn:} 
Đây là phần hướng dẫn chỉ hiện khi bật đáp án.
 \par 


Họ và tên:  ……………… \ Lớp:  ________________ 



ZLZVERBZ00001ZEND



Hoá học: $\ce{2H2 + O2 -> 2H2O}$, đơn vị $9.8\,\mathrm{m}/\mathrm{s}^2$.
\end{document}
